\section{Chapter Preface}
Humans constantly reconstruct hidden reality from fragments. A physician reasons from signs, a physicist from detector traces, and a historian from incomplete documents. These activities look different on the surface but share the same structural core: a projection from latent causes to observed effects.

This chapter universalizes that structure. Reconstruction is presented as a general computational problem independent of disciplinary context, and ambiguity is treated as a geometric property of fibres rather than a defect of a particular model.

The chapter then isolates the conditions under which constraints convert observational equivalence classes into identifiable states.

\section{Derivations and Formal Results}
Let \(\projection:\stateSpace\to\observationSpace\) and define \(\mathcal F_y=\{x\in\stateSpace:\projection(x)=y\}.\)

\begin{theorem}[Local fibre dimension]
If \(\projection:\mathbb R^n\to\mathbb R^m\) is differentiable and \(D\projection_x\) has rank \(r\) on a neighborhood, then locally
\[
\dim \mathcal F_y = n-r.
\]
\end{theorem}

\begin{proof}
This is a direct consequence of the constant-rank theorem: in adapted coordinates, \(\projection\) is locally equivalent to coordinate projection onto \(r\) coordinates, leaving \(n-r\) free coordinates on the fibre.
\end{proof}

\begin{definition}[Reconstruction equivalence]
Define
\[
x\sim_{\projection}x'\iff\projection(x)=\projection(x').
\]
\end{definition}

\begin{proposition}
\(\sim_{\projection}\) is an equivalence relation on \(\stateSpace\).
\end{proposition}

\begin{proof}
Reflexive: \(\projection(x)=\projection(x)\). Symmetric: equality is symmetric. Transitive: if \(\projection(x)=\projection(x')\) and \(\projection(x')=\projection(x'')\), then \(\projection(x)=\projection(x'')\).
\end{proof}

\begin{theorem}[Constraint-induced identifiability]
\label{thm:constraint-induced-identifiability}
Let \(\admissibleSet_{\constraintSet}=\bigcap_i\{x:c_i(x)\le0\}.\) If \(\mathcal F_y\cap\admissibleSet_{\constraintSet}\) contains exactly one element, then the latent state is identifiable relative to \(\constraintSet\).
\end{theorem}
